\begin{corollary}[Endpoint bound for the historical spectral normalization $\mathcal{R} = -\sqrt2\,R_\gamma$]
Let $\mathcal{M}(\gamma_n;\mathbb{R}^n)$ be the space of $\gamma_n$-a.e.\
equivalence classes of measurable $\mathbb{R}^n$-valued functions on
$\mathbb{R}^n$, and let
$T : L^1(\gamma_n;\mathbb{R}) \to \mathcal{M}(\gamma_n;\mathbb{R}^n)$ be any
linear map with the two properties
\begin{enumerate}
\item $T(\iota v) = R_\gamma v$ for every $v \in L^2(\gamma_n)$, with $\iota$ the
inclusion of \noderef{ToL1} and $R_\gamma$ the transform of
\noderef{RieszTransform}; and
\item $t\,\gamma_n\bigl(\{x : |Tf(x)|_{\mathbb{R}^n} > t\}\bigr)
       \leq 2\,\|f\|_{L^1(\gamma_n)}$
for every $f \in L^1(\gamma_n;\mathbb{R})$ and every real $t > 0$.
\end{enumerate}
Then the scalar multiple $(-\sqrt2)\,T$, which is again a linear map
$L^1(\gamma_n;\mathbb{R}) \to \mathcal{M}(\gamma_n;\mathbb{R}^n)$, satisfies
\[
  t\,\gamma_n\bigl(\{x : |((-\sqrt2)\,T)f(x)|_{\mathbb{R}^n} > t\}\bigr)
    \;\leq\; 2\sqrt2\,\|f\|_{L^1(\gamma_n)}
\]
for every $f \in L^1(\gamma_n;\mathbb{R})$ and every real $t > 0$.

This is the second assertion of the reference paper's corollary on coordinate and
historical conventions, its displayed bound
$\sup_{t>0} t\,\gamma_n(\{|\mathcal{R}f|>t\}) \leq 2\sqrt2\|f\|_1$ for the
operator $\mathcal{R} = -\nabla L^{-1/2}\Pi_0$ of the historical
Fabes--Guti\'errez--Scotto normalization, stated per level $t$ rather than as a
supremum.  No new operator has to be introduced for it: the paper's own spectral
computation identifies $\mathcal{R} = -\sqrt2\,R_\gamma$ coordinatewise, the
factor $\sqrt2$ coming from $\partial_j = \sqrt2\,D_j$, so $\mathcal{R}$ is
literally the scalar multiple $(-\sqrt2)\,T$ of the extension $T$.  The content
is therefore a bound for a \emph{different operator} --- one whose weak-type
constant is $2\sqrt2$, not $2$ --- and not a restatement of property (2): the two
inequalities have different left-hand sides and different constants, and neither
constant can be replaced by the other.

Properties (1) and (2) are written out here rather than imported by citation.
They are exactly the two properties for which \noderef{EndpointTheorem} asserts
existence and uniqueness of a map $T$, so the universally quantified statement
above is not vacuous and its instance at that unique $T$ is the paper's
assertion about $\mathcal{R}$.
\end{corollary}
\begin{proof}
Fix $f \in L^1(\gamma_n;\mathbb{R})$ and a real $t > 0$, and write
$\lambda = \sqrt2 > 0$.

\emph{Step 1: the pointwise norm identity.}
Scalar multiplication in $\mathcal{M}(\gamma_n;\mathbb{R}^n)$ is defined
representativewise, and $T$ is linear in the sense that
$((-\lambda)T)f = (-\lambda)\,(Tf)$ as elements of
$\mathcal{M}(\gamma_n;\mathbb{R}^n)$; so for $\gamma_n$-almost every $x$,
\[
  |((-\lambda)T)f(x)|_{\mathbb{R}^n} = |(-\lambda)\,Tf(x)|_{\mathbb{R}^n}
    = |-\lambda|\;|Tf(x)|_{\mathbb{R}^n} = \lambda\,|Tf(x)|_{\mathbb{R}^n},
\]
using absolute homogeneity of the Euclidean norm and $\lambda > 0$.  (The
identity holds almost everywhere rather than at every point because the
displayed objects are a.e.\ equivalence classes; that is all that is used below,
since the sets occurring are measured by $\gamma_n$.)

\emph{Step 2: the level set at $t$ is a level set of $T f$ at $t/\lambda$.}
Let $x$ be a point where the identity of Step 1 holds.  Since $\lambda > 0$,
\[
  t < |((-\lambda)T)f(x)|_{\mathbb{R}^n}
  \iff t < \lambda\,|Tf(x)|_{\mathbb{R}^n}
  \iff \frac{t}{\lambda} < |Tf(x)|_{\mathbb{R}^n},
\]
the last step being division of a strict inequality by the positive number
$\lambda$.  The two sets
$\{x : |((-\lambda)T)f(x)|_{\mathbb{R}^n} > t\}$ and
$\{x : |Tf(x)|_{\mathbb{R}^n} > t/\lambda\}$ therefore differ only inside a
$\gamma_n$-null set, so they have the same $\gamma_n$-measure (monotonicity up to
null sets, shared context, applied in both directions):
\[
  \gamma_n\bigl(\{|((-\lambda)T)f|_{\mathbb{R}^n} > t\}\bigr)
    = \gamma_n\bigl(\{|Tf|_{\mathbb{R}^n} > t/\lambda\}\bigr).
\]

\emph{Step 3: homogeneity of the weak quasi-norm.}
Apply hypothesis (2) to this $f$ at the level $t/\lambda > 0$:
\[
  \frac{t}{\lambda}\,\gamma_n\bigl(\{|Tf|_{\mathbb{R}^n} > t/\lambda\}\bigr)
    \leq 2\,\|f\|_{L^1(\gamma_n)} .
\]
Multiplying by $\lambda > 0$ and using Step 2,
\[
  t\,\gamma_n\bigl(\{|((-\lambda)T)f|_{\mathbb{R}^n} > t\}\bigr)
    = \lambda \cdot \frac{t}{\lambda}\,
      \gamma_n\bigl(\{|Tf|_{\mathbb{R}^n} > t/\lambda\}\bigr)
    \leq \lambda \cdot 2\,\|f\|_{L^1(\gamma_n)}
    = 2\sqrt2\,\|f\|_{L^1(\gamma_n)} ,
\]
which is the assertion.  This is exactly the step the reference paper calls the
homogeneity of the weak-$L^1$ quasi-norm: rescaling the operator by $\lambda$
rescales the level at which the endpoint bound is applied by $1/\lambda$ and the
resulting constant by $\lambda$.
\end{proof}
